\subsection{EMBEDDING OF A MANIFOLD IN THE NEIGHBOURHOOD OF A COMPACT SET}

Lemma 1. Let T and T' be two topological spaces, A and A' compact subsets of T and T', respectively, W a neighbourhood of A × A' in T × T'. There exists an open neighbourhood U of A in T and an open neighbourhood U' of A' in T' such that U × U' ⊂ W.

Let \(x \in \mathrm{A}\) ; there exist open subsets \(\mathrm{U}_x\) of \(\mathrm{T}\) and \(\mathrm{U}_x'\) of \(\mathrm{T}'\) such that \(\{x\} \times \mathrm{A}' \subset \mathrm{U}_x \times \mathrm{U}_x' \subset \mathrm{W}\) : indeed, the compact subset \(\{x\} \times \mathrm{A}'\) of \(\mathrm{T} \times \mathrm{T}'\) can be covered by a finite number of open sets contained in \(\mathrm{W}\) , of the form \(\mathrm{U}_i \times \mathrm{U}_i'\) , with \(x \in \mathrm{U}_i\) ; it suffices to take \(\mathrm{U}_x = \bigcap_i \mathrm{U}_i\) and \(\mathrm{U}_x' = \bigcup_i \mathrm{U}_i'\) .

Since A is compact, there exist points \(x_{1},\ldots,x_{m}\) of A such that \(A\subset\bigcup_{i}U_{x_{i}}\) ; put \(U=\bigcup_{i}U_{x_{i}}\) and \(U'=\bigcap_{i}U'_{x_{i}}\) . Then \(A\times A'\subset U\times U'\subset W\) , hence the lemma.

In the remainder of this number, Y denotes a separated manifold of class \(C^{r}\) .

PROPOSITION 1. Let \(\varphi : \mathrm{X} \to \mathrm{Y}\) be a morphism of class \(\mathbf{C}^r\) , A a compact subset of \(\mathrm{X}\) . The following conditions are equivalent:

\begin{enumerate}
\def\labelenumi{(\roman{enumi})}
\item
  The restriction of \(\varphi\) to A is injective, and \(\varphi\) is an immersion at every point of A;
\item
  there exists an open neighbourhood U of A such that \(\varphi\) induces an embedding of U into Y.
\end{enumerate}

When these conditions are satisfied, \(\varphi\) is said to be an embedding in the neighbourhood of A.

We prove that (i) implies (ii), the other implication being clear. Assuming (i), there exists an open neighbourhood V of X containing A such that the restriction of \(\varphi\) to V is an immersion (Differentiable and Analytic Manifolds, Results, 5.7.1). Denote by \(\Gamma\) the set of points \((x,y)\) in \(V \times V\) such that \(\varphi(x) = \varphi(y)\) , and by \(\Delta\) the diagonal in \(V \times V\) . Then \(\Delta\) is an open subset of \(\Gamma\) : indeed, for all \(x \in V\) , there exists an open neighbourhood \(U_x\) of x such that the restriction of \(\varphi\) to \(U_x\) is injective, that is, such that \(\Gamma \cap (\mathrm{U}_x \times \mathrm{U}_x) = \Delta \cap (\mathrm{U}_x \times \mathrm{U}_x)\) .

Since Y is separated, \(\varGamma\) is closed in \(V \times V\) ; consequently the complement W of \(\varGamma - \varDelta\) in \(V \times V\) is open. By assumption, W contains \(A \times A\) ; Lemma 1 implies that there exists an open subset \(U'\) of V containing A such that \(U' \times U' \subset W\) , that is, such that the restriction of \(\varphi\) to \(U'\) is injective. Moreover, there exists an open neighbourhood U of A whose closure is compact and contained in \(U'\) (General Topology, Chap. I, §9, no. 7, Prop. 10). Then \(\varphi\) induces a homeomorphism from \(\overline{U}\) to \(\varphi(\overline{U})\) , and consequently from U to \(\varphi(U)\) ; it follows that the restriction of \(\varphi\) to U is an embedding (Differentiable and Analytic Manifolds, Results, 5.8.3).

PROPOSITION 2. Assume that the manifold Y is paracompact; let A be a subset of X, and let \(\varphi : X \to Y\) be a morphism of class \(C^r\) that induces a homeomorphism from A to \(\varphi(A)\) , and that is étale at every point of A. Then there exists an open neighbourhood U of A such that \(\varphi\) induces an isomorphism from U to an open submanifold of Y.

Restricting X and Y if necessary, we can assume that \(\varphi\) is étale and surjective. Denote by \(\sigma : \varphi(A) \to A\) the inverse homeomorphism of \(\varphi|A\) . Since Y is metrizable (Differentiable and Analytic Manifolds, Results, 5.1.6), \(\varphi(A)\) admits a fundamental system of paracompact neighbourhoods; hence, by General Topology, Chap. XI, there exist an open neighbourhood V of \(\varphi(A)\) in Y and a continuous map \(s : V \to X\) , that coincides with \(\sigma\) on \(\varphi(A)\) and is such that \(\varphi(s(y)) = y\) for all \(y \in V\) . Moreover, s is topologically étale, so \(s(V)\) is an open set U containing A. Then \(\varphi\) induces a homeomorphism \(\varphi'\) from U to V; by Differentiable and Analytic Manifolds, Results, 5.7.8, \(\varphi'\) is an isomorphism.

In the remainder of this number, we assume that \(r \neq \omega\) .

PROPOSITION 3. Let A be a compact subset of X. The set P of morphisms \(\varphi \in \mathcal{C}^{r}(X;Y)\) that are embeddings in the neighbourhood of A is open in \(\mathcal{C}^{r}(X;Y)\) in the topology of compact \(C^{r}\) -convergence ( \(\S\) 6, no. 4).

Clearly, it suffices to prove the proposition for \(r = 1\) .

\begin{enumerate}
\def\labelenumi{\alph{enumi})}
\tightlist
\item
  We show first that the subset J of \(\mathcal{C}^{1}(\mathrm{X};\mathrm{Y})\) consisting of the morphisms that are immersions at every point of A is open. Consider the map \(j_{\mathrm{A}}:\mathcal{C}^{1}(\mathrm{X};\mathrm{Y})\times\mathrm{A}\to\mathrm{J}^{1}(\mathrm{X};\mathrm{Y})\) such that \(j_{\mathrm{A}}(\varphi,x)=j_{x}^{1}(\varphi)\) (Differentiable and Analytic Manifolds, Results, 12.1).
\end{enumerate}

By definition of the topology on \(\mathcal{C}^1 (\mathrm{X};\mathrm{Y})\) , the map \(\tilde{j}_{\mathrm{A}}:\varphi \to j_{\mathrm{A}}(\varphi ,.)\) from \(\mathcal{C}^1 (\mathrm{X};\mathrm{Y})\) to \(\mathcal{C}(\mathrm{A};\mathrm{J}^1 (\mathrm{X};\mathrm{Y}))\) is continuous; it now follows from General Topology, Chap. X, §3, no. 4, Th. 3, that \(j_{\mathrm{A}}\) is continuous.

On the other hand, let M be the set of jets j in \(\mathrm{J}^{1}(\mathrm{X};\mathrm{Y})\) whose tangent map \(\mathrm{T}(j):\mathrm{T}_{\mathbf{s}(j)}(\mathrm{X})\to\mathrm{T}_{\mathbf{b}(j)}(\mathrm{Y})\) (Differentiable and Analytic Manifolds, Results, 12.3.4) is injective. The set M is open in \(\mathrm{J}^{1}(\mathrm{X};\mathrm{Y})\) ; indeed, it suffices to verify this assertion when X is an open subset of a finite dimensional vector space E, and Y is an open subset of a Banach space F; we are then reduced (Differentiable and Analytic Manifolds, Results, 12.3.1) to proving that the set of injective continuous linear maps is open in \(\mathcal{L}(\mathrm{E};\mathrm{F})\) , which follows from Spectral Theory, Chap. III, §2, no. 7, Prop. 16.

We conclude from the preceding that the set \(j_{\mathrm{A}}^{-1}(\mathrm{M})\) is open in \(\mathcal{C}^{1}(\mathrm{X};\mathrm{Y})\times\mathrm{A}\) , hence that its complement F is closed. Since A is compact, the projection \(\mathrm{pr}_{1}:\mathcal{C}^{1}(\mathrm{X};\mathrm{Y})\times\mathrm{A}\to\mathcal{C}^{1}(\mathrm{X};\mathrm{Y})\) is a proper morphism, hence closed; consequently, the set J, which is equal to \(\mathcal{C}^{1}(\mathrm{X};\mathrm{Y})-\mathrm{pr}_{1}(\mathcal{F})\) , is open in \(\mathcal{C}^{1}(\mathrm{X};\mathrm{Y})\) .

\begin{enumerate}
\def\labelenumi{\alph{enumi})}
\setcounter{enumi}{1}
\tightlist
\item
  Let H be the subset of \(J \times A \times A\) consisting of the elements \((f, x, y)\) such that \(f(x) = f(y)\) . It is clear that H contains \(J \times \Delta\) , where \(\Delta\) denotes the diagonal in the product \(A \times A\) ; we show that \(H' = H - (J \times \Delta)\) is closed in \(J \times A \times A\) . Since P is the complement in J of the image of \(H'\) under the proper projection \(pr_{1}: J \times A \times A \to J\) , this will imply the proposition.
\end{enumerate}

The topology of \(\mathcal{C}^{1}(\mathrm{X};\mathrm{Y})\) being finer than the topology of compact convergence, the map \((\varphi,x)\mapsto\varphi(x)\) from \(\mathcal{C}^{1}(\mathrm{X};\mathrm{Y})\times\mathrm{A}\) to Y is continuous (General Topology, Chap. X, §3, no. 4, Cor. 1 of Th. 3); it follows that H is closed in \(J\times A\times A\) . Hence, it suffices to show that \(J\times\Delta\) is open in H, in other words, that for all \(\varphi\in J\) and all \(x\in A\) there exists a neighbourhood \(\Omega\) of \(\varphi\) in J and a neighbourhood B of x in X such that, for any morphism \(\psi\) in \(\Omega\) , the restriction of \(\psi\) to \(A\cap B\) is injective.

Thus, the proposition follows from the following lemma:

Lemma 2. Let x be a point of X, \(\varphi: X \to Y\) a morphism of class \(C^{1}\) that is an immersion at x. There exists a neighbourhood \(\Omega\) of \(\varphi\) in \(\mathcal{C}^{1}(X;Y)\) and a neighbourhood B of x in X such that, for all \(\psi \in \Omega\) , the restriction of \(\psi\) to B is injective.

Let U be a relatively compact open neighbourhood of x isomorphic to a finite dimensional vector space, and such that \(\varphi(\overline{\mathrm{U}})\) is contained in the domain V of a chart. The set \(\Omega_{0}\) of \(\psi\in\mathcal{C}^{1}(\mathrm{X};\mathrm{Y})\) such that \(\psi(\overline{\mathrm{U}})\subset\mathrm{V}\) is open in \(\mathcal{C}^{1}(\mathrm{X};\mathrm{Y})\) , and the restriction map \(\Omega_{0}\to\mathcal{C}^{1}(\mathrm{U};\mathrm{V})\) is continuous; we are thus reduced to proving the lemma when X=U and Y=V, in other words, we can assume that X is a finite dimensional vector space and Y is a Banach space. Choose norms on X and Y.

The linear map \(\mathrm{D}\varphi(x):\mathrm{X}\to\mathrm{Y}\) is injective; denote its conorm by q (Spectral Theory, Chap. III, §2, no. 6), so that, by definition, we have \(\|\mathrm{D}\varphi(x).t\|\geq q\|t\|\) for all \(t\in X\) . Let \(\varepsilon\in R\) be such that \(0<\varepsilon<q/2\) , and let B be a closed ball with centre x such that \(\|\mathrm{D}\varphi(u)-\mathrm{D}\varphi(x)\|\leq\varepsilon\) for all \(u\in B\) . Denote by \(\Omega\) the subset of \(\mathcal{C}^{1}(\mathrm{X};\mathrm{Y})\) consisting of the morphisms \(\psi\) such that \(\|\mathrm{D}\psi(u)-\mathrm{D}\varphi(u)\|\leq\varepsilon\) for all \(u\in B\) ; it is open by definition of the topology of \(\mathcal{C}^{1}(\mathrm{X};\mathrm{Y})\) . For \(\psi\in\Omega\) , put \(\psi_{0}=\psi-\mathrm{D}\varphi(x)\) . We have \(\|\mathrm{D}\psi_{0}(u)\|\leq2\varepsilon\) for all \(u\in B\) , and consequently \(\|\psi_{0}(u)-\psi_{0}(v)\|\leq2\varepsilon\|u-v\|\) for all u and v in B (Differentiable and Analytic Manifolds, Results, 2.2.3). It follows that

\[
\| \psi (u) - \psi (v) \| \geq \| \mathrm{D} \varphi (x). (u - v) \| - \| \psi_ {0} (u) - \psi_ {0} (v) \| \geq (q - 2 \varepsilon) \| u - v \|.
\]

Consequently, the restriction of \(\psi\) to B is injective, hence the lemma.

PROPOSITION 4. Let A be a compact subset of X. There exist a finite dimensional vector space E and a morphism \(\varphi \in \mathcal{C}^r (\mathrm{X};\mathrm{E})\) ( \(r\neq \omega\) ) that is an embedding in the neighbourhood of A.

Let \((\mathrm{U}_{i},\varphi_{i},\mathrm{E}_{i})_{i\in\mathrm{I}}\) be a finite family of charts of X whose domains cover A. We extend \(\varphi_{i}\) to a map from X to \(E_{i}\) (also denoted by \(\varphi_{i}\) ) by putting \(\varphi_{i}(x)=0\) for \(x\notin U_{i}\) . Let \((\mathrm{V}_{i})_{i\in\mathrm{I}}\) be a covering of A by open subsets of X such that \(\overline{\mathrm{V}}_{i}\subset\mathrm{U}_{i}\) for all \(i\in\mathrm{I}\) (the existence of such a covering follows from General Topology, Chap. IX, §4, no. 3, Cor. 1 of Th. 3, applied to the compact space \(X'\) obtained by adjoining to X a point at infinity and the covering of \(X'\) consisting of the open sets \(U_{i}\) ( \(i\in\mathrm{I}\) ) and \(X'-\mathrm{A}\) ). For all \(i\in\mathrm{I}\) , let \(\alpha_{i}\) be a numerical function of class \(C^{r}\) on X, equal to 1 at every point of \(V_{i}\) , and with support contained in \(U_{i}\) (Differentiable and Analytic Manifolds, Results, 5.3.6).

Consider the map \(\varphi : \mathrm{X} \to \bigoplus_{i \in \mathrm{I}} (\mathrm{E}_i \oplus \mathbf{R})\) defined by

\[
\varphi (x) = (\alpha_ {i} (x) \varphi_ {i} (x), \alpha_ {i} (x)) _ {i \in I}.
\]

For all \(i \in I\) , the map \(\alpha_{i}\varphi_{i}\) is of class \(C^{r}\) (since its restrictions to \(U_{i}\) and to the complement of the support of \(\alpha_{i}\) are), and its restriction to \(V_{i}\) is an embedding; it follows that \(\varphi\) is a morphism of class \(C^{r}\) and is an immersion at every point of A. We show that the restriction of \(\varphi\) to A is injective. Let x, y be two points of A such that \(\varphi(x) = \varphi(y)\) , and let \(i \in I\) be such that \(x \in V_{i}\) . Then \(\alpha_{i}(x) = 1\) , so \(\alpha_{i}(y) = 1\) , which implies that \(y \in U_{i}\) ; but we also have \(\varphi_{i}(x) = \varphi_{i}(y)\) , so x = y since \(\varphi_{i}\) induces an embedding of \(U_{i}\) into \(E_{i}\) .

It can be shown \(^{9}\) that every separated manifold, countable at infinity and of pure dimension n, embeds in \(R^{2n}\) ; for a weaker result, cf.~Exercise 2.

